\section{混合項を検出する境界の一般化Weyl表示}
\label{sec:boundary-notes-v39}

\subsection{正規元の選択を変える原理}
$S_H=B[w;\sigma,x^{-1}[-,H]]$、$J=xw+H$ において、
$x^\ell$ を反転した中心開部分は $H$ に依存しない。
一方 $u^\ell$ を反転すれば
\[
 B[u^{-\ell}]=\mathbb C[Q,u^{\pm1}][x;\tau],\qquad
 \tau(u)=t^{-1}u
\]
となり、$H=h(Q,u)+xG$ と一意分解できる。
$W=w+G$ とするだけで
\[
 xW=J-h(Q,u),\qquad Wx=J-h(Q,tu)
\]
という一般化Weyl表示になる。
この変換は、$x$ が全 $S_H$ の正規元であることを必要としない。
$u^\ell$ は本当に全代数の中心なので、中心の境界を扱う正当な局所化である。

\subsection{境界の記憶は有限ノルムに集約される}
$U=u^\ell$、$\Xi=x^\ell$、$\Omega=W^\ell$ と置くと、
\[
 Z(S_H)[U^{-1}]
 =\mathbb C[Q,J,U^{\pm1},\Xi,\Omega]/(\Xi\Omega-P_h),\qquad
 P_h=\prod_{i=0}^{\ell-1}(J-h(Q,t^iu)).
\]
同じノルムは、自由加群 $\mathbb C[Q,u^{\pm1}]$ の
$\mathbb C[Q,U^{\pm1}]$ 上の乗法行列から $\det(JI-M_h)$ と計算できる。
二方法を独立に比較すると、巡回向き、負冪、ノルムの符号を検出できる。

中心は $D=\mathbb C[Q,J,U^{\pm1},\Xi]$ の部分環 $D[P_h/\Xi]$ である。
従って、混合項はアフィン改変の中心 $(\Xi,P_h)$ を動かす。
$Q,U,J$ を固定した $P_h$ の一致は、体内で積の一つの因子が零になることから、
$h'=h(Q,t^ku)$ と同値である。
この復元性は標識付きであり、$J$ の移動を許す分類とは区別する。
特に中心元を $H$ に加えてもOre導分は変わらず、$J$ が移動するだけである。

\subsection{切断間隔が単純加群の次元を与える}
中心点で $\Xi=\Omega=0$ とすると、軌道上の条件
$J=h(Q,t^i\lambda)$、$\lambda^\ell=U$ の成立箇所が切断辺になる。
その個数を $r$、切断後の鎖の頂点数を $d_1,\ldots,d_r$ とする。
基本代数は両向きの巡回箙で、折返しと一周を零にした $\mathcal R_r$ である。
元のファイバーは $\bigoplus(e_i\mathcal R_r)^{d_i}$ の自己準同型代数であり、
\[
 \text{simple dimensions}=(d_i),\qquad
 \dim=2\ell^2-\sum_i d_i^2,\qquad
 \operatorname{nilindex}(\operatorname{rad})=r.
\]
非零辺を縮約する上界と、行列単位・切断矢を用いた明示表現の下界を両方示す。
単純加群の数だけからファイバー全体の次元を推測しない。
$r=1$ の近傍で行列環表示を作ることで、Azumaya性に必要な局所自由性も得られる。
この中心開部分ではAzumaya部分と滑らかな部分が一致する。

\subsection{一般の巡回モノドロミー}
$h$ の巡回安定化群の位数を $g$、$n=\ell/g$ とする。
被約スペクトル因子の正規化は $\mathbb C[Q,u^{\pm g}]$ であり、
$P_h=P_{\mathrm{min}}^g$ となる。
$z=u^g$、$h_i=h(Q,t^iu)$、
$\Delta=\prod_{i=1}^{n-1}(h_0-h_i)$ を用い、追加衝突を除く部分
$A=\mathbb C[Q,z^{\pm1},\Delta^{-1}]$ を取る。
$L=A[\kappa]/(\kappa^g-z)$ 上の角は
\[
 \mathcal T_g=
 L\langle P,R\rangle/
 (Pa-\sigma(a)P,\ Ra-\sigma^{-1}(a)R,\ PR,\ RP,\ P^g,\ R^g),
 \quad \sigma(\kappa)=t^{-n}\kappa.
\]
正規化上の全代数は $M_n(\mathcal T_g)$ である。
その冪零部分を除く商 $M_n(L)$ が、連結な巡回 $g$ 重被覆上の
実際の $n$ 次元単純表現族を与える。
これは任意の $g$ に対する構成であり、追加衝突点では切断間隔の一般則へ戻る。
$v38$ の偶数位数の二単純加群は、この仕組みの二次の場合に対応する。

\subsection{比較に使える最小例と残る問題}
重み $m=4$、冪根位数 $\ell=5$ で、
$H=u^5$ と $H=u^4v$ はそれぞれ
\[
 P_h=(J-U)^5,\qquad P_h=J^5+Q^5U^3
\]
を与える。$U\ne0$ の中心の特異部分の次元は二と一で異なる。
前者の一般境界ファイバーは五つの一次元単純加群を持ち、
後者は $Q\ne0$ で五次元単純加群を一つ持つ。
両者の関数斜体と $\Xi\ne0$ の中心開部分は同じである。
この例は、関数斜体による分類が混合項を検出しないことの具体的な補完となる。

混合項が中心元と $xG$ の和となる場合の全中心と深い境界は、後続節で解決する。
一般の混合項について残る問題は、$\Xi=U=V=0$ の深い境界における全中心と表現、
三つの中心開部分を含む大域的な延長、標識を忘れた座標環・射影圏の分類である。
今回のノルム復元だけから、全 $H$ に対する導来非同値は結論しない。

検証コードは \texttt{verify\_endpoint\_boundary\_v39.py}。
位数 $3,4,5,6,8$ の非可換ゲージ変換と独立な行列式、
位数三から八までの全 $498$ 個の非空切断配置、
および $60$ 個の単項式ノルムを有理数・円分体上で確認する。
一般巡回降下では $32$ 個の $(\ell,g)$ の組について、
$411$ 個の行列恒等式も厳密に検証する。
全位数の主張は原稿中の一般証明による。
